\\documentclass[11pt]{article}
\\usepackage[margin=1in]{geometry}
\\usepackage{booktabs}
\\usepackage{amsmath}
\\usepackage{hyperref}
\\title{MARS: Uncertainty-Aware Adaptive Reasoning for Selective LLM Routing}
\\author{Garrett Carroll \\ Independent Researcher}
\\date{\\today}

\\begin{document}
\\maketitle

\\begin{abstract}
This manuscript presents the evaluation protocol for MARS, an uncertainty-aware routing framework for language-model workflows. MARS records an uncertainty estimate and may allocate additional inference-time reasoning to cases above a pre-registered threshold. The paper makes no empirical performance claim until the corresponding reproducible evaluation artifacts are available. We evaluate accuracy, calibration, error detection, selective risk, latency, and cost against fixed-budget baselines.
\\end{abstract}

\\section{Introduction}
Language-model systems can be accurate on average while providing unreliable confidence on individual outputs. MARS investigates whether a calibrated uncertainty signal can support selective escalation to a deeper reasoning path. The contribution is an explicitly auditable experimental design, not a claim that uncertainty-aware routing is categorically novel.

\\section{Method}
For each input $x$, the system emits confidence $c(x)\\in[0,1]$ and uncertainty $u(x)=1-c(x)$. Under the adaptive condition, a deep route is selected when $u(x)\\geq\\tau$, where $\\tau$ is fixed before held-out testing. All conditions share the same task items, verifier, model version, and prompt version except for the routing intervention.

\\section{Experimental Design}
We compare: (1) a one-pass baseline, (2) self-reported confidence, (3) MARS score-only, (4) MARS adaptive routing, and (5) an always-deep condition. Primary metrics are accuracy, Expected Calibration Error, Brier score, error-detection AUROC, risk-coverage, cost, latency, and escalation rate.

\\section{Results}
\\textbf{Pending reproducible experiment execution.} No numerical result should appear in this section until linked to versioned prediction records, an evaluation manifest, and the scoring code revision.

\\begin{table}[h]
\\centering
\\caption{Results template. Values are intentionally omitted pending execution.}
\\begin{tabular}{lrrrrrr}
\\toprule
Condition & Accuracy & ECE & Brier & AUROC & Cost/Correct & Latency \\\\
\\midrule
Baseline & TBD & TBD & TBD & TBD & TBD & TBD \\\\
Self-reported confidence & TBD & TBD & TBD & TBD & TBD & TBD \\\\
MARS score-only & TBD & TBD & TBD & TBD & TBD & TBD \\\\
MARS adaptive & TBD & TBD & TBD & TBD & TBD & TBD \\\\
Always-deep & TBD & TBD & TBD & TBD & TBD & TBD \\\\
\\bottomrule
\\end{tabular}
\\end{table}

\\section{Limitations}
The system may depend on the reliability of the underlying model, the task verifier, and the calibration dataset. Routing can increase cost or latency, and a confidence signal can fail under distribution shift. Results must be reported by task category rather than only as a pooled average.

\\section{Reproducibility and Ethics}
Every run records model identifiers, prompts, seeds, routing threshold, timestamps, software environment, token use, latency, and costs. Customer or proprietary data must not be included without a documented legal basis and de-identification review.

\\bibliographystyle{plain}
\\bibliography{references}
\\end{document}
